\section{Primera parte: incorporar el pensamiento informal a la demostración formal — LeanSTaR}

\subsection{Motivación: más allá del entrenamiento puramente formal}

Los métodos tradicionales entrenan el modelo solo con datos de demostración formal (estado $\to$ política del siguiente paso), pero esto puede no bastar para aprender el \textbf{proceso de pensamiento} necesario para producir el código.

\begin{importantbox}{La idea central de LeanSTaR}
Entrenar al modelo para que, antes de cada paso de demostración formal, primero genere un \textbf{pensamiento informal} (informal thought) de forma libre:
\begin{itemize}
    \item El modelo puede descomponer el problema, planear el siguiente paso y ejecutar cálculos informales
    \item El pensamiento ayuda a diversificar el espacio de búsqueda
    \item De forma similar a Chain-of-Thought, aumenta la capacidad expresiva del modelo
\end{itemize}
\end{importantbox}

\subsection{El algoritmo de LeanSTaR}

Entrenamiento en dos etapas:

\textbf{Etapa uno: inicialización}
\begin{enumerate}
    \item A partir de demostraciones existentes (como Mathlib), dado el estado y la política del siguiente paso, se le pide al LLM que genere retrospectivamente el pensamiento que describe esa transición
    \item Se recopilan cientos de miles de muestras (estado, pensamiento, política)
    \item Se entrena un modelo inicial para que aprenda a generar el pensamiento antes de la política
\end{enumerate}

\textbf{Etapa dos: aprendizaje por refuerzo (Expert Iteration)}
\begin{enumerate}
    \item El modelo genera en paralelo múltiples demostraciones completas
    \item Lean verifica si la demostración es correcta
    \item Las trayectorias exitosas se recopilan y se agregan al conjunto de entrenamiento
    \item Se entrena un nuevo modelo sobre el conjunto de datos ampliado
    \item Se itera este proceso
\end{enumerate}

\begin{warningbox}{La dificultad de la búsqueda en árbol al incorporar el pensamiento}
La Best-First Search tradicional funciona mal una vez que se incorpora el pensamiento informal: es difícil puntuar el pensamiento con precisión. LeanSTaR opta en cambio por una estrategia sencilla de \textbf{muestreo en paralelo}: se generan varias rutas en paralelo, se reintenta cuando hay un error y se detiene en cuanto una tiene éxito.
\end{warningbox}

\subsection{Resultados experimentales}

\begin{itemize}
    \item En el punto de referencia mini-F2F, LeanSTaR llevó al modelo de estar por debajo de la referencia del agente GPT-4 a superarla
    \item \textbf{Hallazgo clave}: el modelo con pensamiento se beneficia más de la escalabilidad del presupuesto de inferencia (scaling inference budget), porque el pensamiento diversifica de manera efectiva el espacio de búsqueda
    \item En el pensamiento informal, el modelo ejecutó pasos de razonamiento que no existen en el código formal
\end{itemize}

Desde entonces, el pensamiento informal combinado con entrenamiento por RL ha sido adoptado ampliamente por la comunidad de demostración de teoremas (DeepSeek-Prover, entre otros), y también es el paradigma central de modelos de razonamiento como O1 y DeepSeek-R1.

\subsection{Resumen de la sección}

La idea central de LeanSTaR: entrenar solo con código formal puede no bastar para aprender el proceso de pensamiento necesario para generar código. Aprender a generar pensamiento informal mediante RL puede mejorar de manera significativa la capacidad de demostración.
